% ===== PRÉAMBULE CANONIQUE — à coller VERBATIM en tête de CHAQUE .tex =====
\documentclass[11pt,a4paper]{article}
\usepackage[utf8]{inputenc}\usepackage[T1]{fontenc}\usepackage{lmodern}
\usepackage[french]{babel}
\usepackage{amsmath,amssymb,mathtools}\usepackage{icomma}
\usepackage{siunitx}\sisetup{locale=FR}  % \qty{}{}, \si{}, \ang{}
\usepackage{xcolor}\usepackage{colortbl}
\usepackage{tikz}\usepackage{pgfplots}\pgfplotsset{compat=1.18}
\usepackage[siunitx,europeanresistors]{circuitikz} % circuits (élec.)
\usepackage[most]{tcolorbox}\usepackage{enumitem}\usepackage{array,booktabs}
\usepackage[a4paper,margin=2.2cm]{geometry}
\usetikzlibrary{arrows.meta,calc,angles,quotes,positioning,patterns,%
  backgrounds,shapes.geometric,decorations.pathmorphing,decorations.markings,intersections}
\definecolor{primary}{RGB}{222,49,99}\definecolor{primarydark}{RGB}{150,22,55}
\definecolor{defcol}{RGB}{0,114,178}\definecolor{methcol}{RGB}{52,92,148}
\definecolor{excol}{RGB}{0,135,90}\definecolor{attncol}{RGB}{200,40,55}
\tcbset{boxbase/.style={enhanced,breakable,boxrule=0.9pt,arc=2.5pt,
  left=3mm,right=3mm,top=2.3mm,bottom=2.3mm,fonttitle=\bfseries,coltitle=white}}
% --- boîtes TUTORIEL (noms EXACTS, ne pas renommer) ---
\newtcolorbox{cle}[1][]{boxbase,colback=defcol!6,colframe=defcol,
  title={\textsc{À retenir}\ifx\relax#1\relax\relax\else\textmd{ : \itshape #1}\fi}}
\newtcolorbox{methode}[1][]{boxbase,colback=methcol!7,colframe=methcol,sharp corners,
  title={\textsc{Méthode}\ifx\relax#1\relax\relax\else\textmd{ --- \itshape #1}\fi}}
\newtcolorbox{exemplebox}[1][]{boxbase,colback=excol!5,colframe=excol,
  title={\textsc{Exemple}\ifx\relax#1\relax\relax\else\textmd{ : \itshape #1}\fi}}
\newtcolorbox{piege}[1][]{boxbase,colback=attncol!6,colframe=attncol,
  title={\textsc{Piège}\ifx\relax#1\relax\relax\else\textmd{ : \itshape #1}\fi}}
\newtcolorbox{remarque}[1][]{boxbase,colback=black!4,colframe=black!45,
  title={\textsc{Remarque}\ifx\relax#1\relax\relax\else\textmd{ : \itshape #1}\fi}}
% --- boîtes EXERCICE / CORRIGÉ (noms EXACTS, auto-comptées) ---
\newtcolorbox[auto counter]{exo}[1][]{boxbase,colback=primary!4,colframe=primary,
  title={\textsc{Exercice}~\thetcbcounter\ifx\relax#1\relax\relax\else\textmd{ --- \itshape #1}\fi}}
\newtcolorbox[auto counter]{corrige}[1][]{boxbase,colback=excol!4,colframe=excol,
  title={\textsc{Corrigé}~\thetcbcounter\ifx\relax#1\relax\relax\else\textmd{ --- \itshape #1}\fi}}
\newcommand{\vv}[1]{\vec{#1}}\newcommand{\ds}{\displaystyle}
\newcommand{\R}{\mathbb{R}}\newcommand{\N}{\mathbb{N}}\newcommand{\Z}{\mathbb{Z}}
% =========================================================================
\begin{document}
\sisetup{per-mode=symbol}

\begin{corrige}[le recul du canon, en chiffres]
\begin{enumerate}[label=\alph*)]
  \item $a_{\text{boulet}}=\dfrac{F}{m_{\text{boulet}}}=\dfrac{4000}{8}=\qty{500}{\metre\per\second\squared}$.
  \item $a_{\text{canon}}=\dfrac{F}{m_{\text{canon}}}=\dfrac{4000}{800}=\qty{5}{\metre\per\second\squared}$ (sens opposé).
  \item $\dfrac{a_{\text{boulet}}}{a_{\text{canon}}}=\dfrac{500}{5}=100=\dfrac{800}{8}=\dfrac{m_{\text{canon}}}{m_{\text{boulet}}}$. \checkmark
\end{enumerate}
\end{corrige}

\begin{corrige}[la feuille de papier frappée]
\begin{enumerate}[label=\alph*)]
  \item $a=\dfrac{\Delta v}{\Delta t}=\dfrac{30-0}{\num{0,06}}=\qty{500}{\metre\per\second\squared}$ ; avec
        $m=\qty{5}{\gram}=\qty{0,005}{\kilogram}$ : $F=ma=\num{0,005}\times500=\qty{2,5}{\newton}$.
  \item Par la 3\textsuperscript{e} loi, la feuille exerce sur la main une force de \emph{même intensité} :
        $\qty{2,5}{\newton}$ (impossible d'exercer plus que ce qu'on reçoit).
\end{enumerate}
\end{corrige}

\begin{corrige}[la propulsion d'une fusée]
\begin{enumerate}[label=\alph*)]
  \item $a_{\text{gaz}}=\dfrac{F}{m_{\text{gaz}}}=\dfrac{8000}{4}=\qty{2000}{\metre\per\second\squared}$ (vers le bas).
  \item Le gaz repousse la fusée avec la même force (\qty{8000}{\newton}) vers le haut :
        $a_{\text{fusée}}=\dfrac{8000}{200}=\qty{40}{\metre\per\second\squared}$, dirigée \emph{vers le haut}.
\end{enumerate}
\end{corrige}

\begin{corrige}[la moto et la Terre]
\begin{enumerate}[label=\alph*)]
  \item $a_{\text{moto}}=\dfrac{F}{m}=\dfrac{800}{200}=\qty{4}{\metre\per\second\squared}$.
  \item $a_{\text{Terre}}=\dfrac{F}{M}=\dfrac{800}{\num{6e24}}\approx\qty{1,3e-22}{\metre\per\second\squared}$ :
        absolument imperceptible à cause de l'énorme masse de la Terre (inertie colossale).
\end{enumerate}
\end{corrige}

\begin{corrige}[sauter d'une barque]
\begin{enumerate}[label=\alph*)]
  \item $a_1=\dfrac{F}{m_1}=\dfrac{180}{60}=\qty{3}{\metre\per\second\squared}$ ;
        $a_2=\dfrac{F}{m_2}=\dfrac{180}{120}=\qty{1,5}{\metre\per\second\squared}$.
  \item La barque part \emph{en sens opposé} à la personne (recul). La personne, plus légère, accélère le plus.
\end{enumerate}
\end{corrige}
\end{document}
